\begin{theorem}[uniform non-torsion]
\label{thm:uniform-nontorsion-audit}
\Verified\ There is a line bundle \(L\in\operatorname{Pic}(X)\) such that
\(L|_{X_b}\) is non-torsion for every scheme-theoretic fibre
\(X_b=f^*(b)\), and \(L|_{S_j}\) is non-torsion for \(j=1,2\).
More precisely, identify \(\operatorname{Pic}(X)\) with \(\mathbb C\),
\(z\mapsto L_z\).  Every \(z\) outside the union of \(\mathbb R\) with a
countable set has these properties.
\end{theorem}

\begin{proof}
The pullback \(f^*\mathcal O_B(1)\) is nontrivial: by the projection
formula and \(f_*\mathcal O_X=\mathcal O_B\),
\[
 h^0(X,f^*\mathcal O_B(1))=h^0(B,\mathcal O_B(1))=2,
\]
whereas \(h^0(X,\mathcal O_X)=1\).  Rescale the additive coordinate in
\cref{cor:picard-line-audit} so that
\[
 L_{12}=f^*\mathcal O_B(1).
\tag{9.19}
\]

Fix a smooth fibre \(F_q\).  Put
\[
 \Lambda_q=\iota_q^{-1}\!\left(
 \operatorname{im}[H^1(F_q;\mathbb Z)\to
 H^1(F_q,\mathcal O_{F_q})]\right)\subset\mathbb C.
\tag{9.20}
\]
The integral image is a lattice in the underlying real vector space of
\(H^1(F_q,\mathcal O_{F_q})\), so \(\Lambda_q\) is discrete.  It contains
\(12\), because (9.19) is trivial on every fibre.  The exponential
sequence on \(F_q\) shows
\[
 L_z|_{F_q}\text{ is torsion}
 \quad\Longleftrightarrow\quad z\in\mathbb Q\Lambda_q.
\tag{9.21}
\]
If \(\operatorname{rk}\Lambda_q=1\), its nonzero element \(12\) forces
\(\Lambda_q\subset\mathbb R\); all its bad parameters are real.

Suppose \(\operatorname{rk}\Lambda_q=2\).  By
\cref{cor:picard-line-audit}, the image of \(\iota_q\) is the line
generated by the image of \(\gamma\).  Thus there is an integral class
\[
 \lambda=a\gamma+bu+cw+d\delta\notin\mathbb Q\gamma
\tag{9.22}
\]
whose \((0,1)\)-image lies on that line.  On the covering family, the four
marked lattice columns are
\[
 \pi_{\widehat\gamma}=(6\mu,\beta),\quad
 \pi_{\widehat u}=(\tau,\mu),\quad
 \pi_{\widehat w}=(1,0),\quad
 \pi_{\widehat\delta}=(0,1).
\tag{9.23}
\]
Every real cohomology class has a unique expression
\[
 2\operatorname{Re}(\bar\xi\,d\zeta_1+
                       \bar\upsilon\,d\zeta_2).
\]
Evaluating it on (9.23) gives
\[
\begin{aligned}
2\operatorname{Re}\bar\xi&=c,&
2\operatorname{Re}\bar\upsilon&=d,\\
2\operatorname{Re}(\bar\xi\tau+\bar\upsilon\mu)&=b,&
2\operatorname{Re}(6\bar\xi\mu+\bar\upsilon\beta)&=a.
\end{aligned}
\tag{9.24}
\]
For \(\gamma\), the first two right sides vanish.  The third equation then
shows that its \((0,1)\)-line is represented, up to a nonzero scalar, by
\((\operatorname{Im}\mu,-\operatorname{Im}\tau)\); the fourth equation
only fixes the scalar.  Write
\(\bar\xi=c/2+ir\), \(\bar\upsilon=d/2+is\).  The condition that the class
in (9.22) have \((0,1)\)-image on the \(\gamma\)-line is
\[
 \bar\xi\operatorname{Im}\tau+
 \bar\upsilon\operatorname{Im}\mu=0.
\]
Its real part, together with the third equation of (9.24), is equivalent
to the single holomorphic equation
\[
 c\tau+d\mu-b=0.
\tag{9.25}
\]
Neither \(a\) nor \(\beta\) occurs, as adding a multiple of \(\gamma\)
must not change the condition.

The functions \(1,\tau,\mu\) are linearly independent.  Applying the cusp
transformation \((\tau,\mu)\mapsto(\tau-1,\mu)\) to a constant linear
relation first kills the coefficient of \(\tau\).  If the remaining
relation made \(\mu\) constant, the order-four law
\(\mu\circ g_2=1+\mu/\tau\), with nonconstant \(\tau\), would force an
impossible identity.  Hence all coefficients vanish.

For every integral \(\lambda\notin\mathbb Q\gamma\), equation (9.25) is
therefore a nonzero holomorphic function on the covering Riemann surface.
Its zero set is discrete and countable.  There are countably many
\(\lambda\), so their union, and its image in \(B^\circ\), are countable.
This image contains every \(q\) for which
\(\operatorname{rk}\Lambda_q=2\).  Each such fibre contributes the
countable bad set \(\mathbb Q\Lambda_q\); all rank-one fibres together
contribute only parameters in \(\mathbb R\).  Thus the smooth fibres
exclude only \(\mathbb R\) together with a countable set.

There remain the three scheme-theoretic special fibres.  By
\cref{thm:structure-direct-images-audit,cor:picard-line-audit}, restriction
\[
 \mathbb C\longrightarrow H^1(X_b,\mathcal O_{X_b})
\]
is complex-linear and injective for each of them.  The exponential sequence
and countability of \(H^1(X_b;\mathbb Z)\) show that the parameters for
which \(L_z|_{X_b}\) is torsion form a countable set.

Finally consider the reductions \(S_j\).  The restriction
\(\varphi_j:\mathbb C\to H^1(S_j,\mathcal O_{S_j})\) is complex-linear.
It is nonzero.  Write \(\mathcal O_X(S_j)=L_{z_j}\).  Since
\[
 \mathcal O_X(m_jS_j)=f^*\mathcal O_B(1)=L_{12},
\]
the additive parameter gives \(z_j=12/m_j\), while
\cref{lem:finite-normal-order-audit} says that
\[
 L_{12/m_j}|_{S_j}=\mathcal O_X(S_j)|_{S_j}
\]
is nontrivial.  If \(\varphi_j\) were zero, the exponential sequence and
the vanishing of the restricted first Chern class would make every
\(L_z|_{S_j}\) trivial, a contradiction.  Thus \(\varphi_j\) is
injective, and the same exponential-sequence argument makes its torsion
parameters countable.  Removing these two additional countable sets proves
the theorem.
\end{proof}
